% GENERATED FILE. Edit canonical lessons, then run npm run generate:books.
% canonical-source-hash: fnv1a64:efc992badd6b3ff1
% canonical-lessons: PT-C504-colecao, PT-C504-colecionar, PT-C504-artesanato, PT-C504-barro, PT-C504-olaria

\chapter{Coleções e artesanato}
\label{ch:pt-colecoes-e-artesanato}
\hlchaptermodality{504}

\begin{chapteropening}
\textbf{By the end of this chapter:} \emph{I can talk about collecting and crafts.}

Complete the last lesson of chapter 504: I can talk about collecting and crafts.
\end{chapteropening}

\section[a coleção]{a coleção — collection}

\label{lesson:PT-C504-colecao}



\begin{quote}
Before the new one: say the Portuguese for jigsaw puzzle, then the Portuguese for crossword.
\end{quote}

\subsection*{You'll want to know: a coleção}
\textbf{a coleção} — \textquotedblleft{}collection\textquotedblright{}.

\begin{quote}
Tem uma coleção enorme de discos.
\end{quote}

\begin{quote}
\emph{He has a huge record collection.}
\end{quote}

Before AO90 spelled \emph{colecção} in Portugal.

\begin{cousinweb}
\emph{coleção} is a learned borrowing from Latin \emph{collēctiō} \textquotedblleft{}a gathering together\textquotedblright{} (EtymDB), from \emph{colligere} \textquotedblleft{}to gather\textquotedblright{}.
\end{cousinweb}

\begin{grammarlens}[title={What you've built}]
One more word for telling, explaining and describing at B1.
\end{grammarlens}

\subsection*{Your turn}
\begin{itemize}
\raggedright
  \item \emph{Say it:} \emph{a coleção}
  \item \emph{Say it:} \emph{a coleção}, once more
  \item \emph{Say it:} a sentence of your own with \emph{a coleção}
  \item \emph{Recall:} say the Portuguese for jigsaw puzzle, then the Portuguese for crossword, then say \emph{a coleção} again
\end{itemize}

\begin{tcolorbox}[breakable,colback=teal!4,colframe=teal!35!black,title={Before you move on}]
What does a coleção mean? (\textquotedblleft{}Collection\textquotedblright{}.) Say it once more.
\end{tcolorbox}
\section[colecionar]{colecionar — to collect}

\label{lesson:PT-C504-colecionar}



\begin{quote}
Before the new one: say the Portuguese for crossword, then the Portuguese for collection.
\end{quote}

\subsection*{You'll want to know: colecionar}
\textbf{colecionar} — \textquotedblleft{}to collect\textquotedblright{}.

\begin{quote}
O meu filho coleciona cromos de futebol.
\end{quote}

\begin{quote}
\emph{My son collects football stickers.}
\end{quote}

Before AO90 spelled \emph{coleccionar}. Built on \emph{a coleção}.

\begin{cousinweb}
\emph{colecionar} is built on \emph{coleção} \textquotedblleft{}collection\textquotedblright{}, a learned borrowing from Latin \emph{collēctiō} (EtymDB).
\end{cousinweb}

\begin{grammarlens}[title={What you've built}]
One more word for telling, explaining and describing at B1.
\end{grammarlens}

\subsection*{Your turn}
\begin{itemize}
\raggedright
  \item \emph{Say it:} \emph{colecionar}
  \item \emph{Say it:} \emph{colecionar}, once more
  \item \emph{Say it:} a sentence of your own with \emph{colecionar}
  \item \emph{Recall:} say the Portuguese for crossword, then the Portuguese for collection, then say \emph{colecionar} again
\end{itemize}

\begin{tcolorbox}[breakable,colback=teal!4,colframe=teal!35!black,title={Before you move on}]
What does colecionar mean? (\textquotedblleft{}To collect\textquotedblright{}.) Say it once more.
\end{tcolorbox}
\section[o artesanato]{o artesanato — handicrafts, crafts}

\label{lesson:PT-C504-artesanato}



\begin{quote}
Before the new one: say the Portuguese for collection, then the Portuguese for to collect.
\end{quote}

\subsection*{You'll want to know: o artesanato}
\textbf{o artesanato} — \textquotedblleft{}handicrafts, crafts\textquotedblright{}.

\begin{quote}
Na feira comprámos artesanato da região.
\end{quote}

\begin{quote}
\emph{At the fair we bought local crafts.}
\end{quote}

The craftsman is \emph{o artesão}.

\begin{cousinweb}
\emph{artesanato} is built on \emph{artesão} \textquotedblleft{}craftsman\textquotedblright{}, which is usually traced to Italian \emph{artigiano}; that word goes back to Latin \emph{ars} \textquotedblleft{}skill, craft\textquotedblright{} (Lewis \& Short).
\end{cousinweb}

\begin{grammarlens}[title={What you've built}]
One more word for telling, explaining and describing at B1.
\end{grammarlens}

\subsection*{Your turn}
\begin{itemize}
\raggedright
  \item \emph{Say it:} \emph{o artesanato}
  \item \emph{Say it:} \emph{o artesanato}, once more
  \item \emph{Say it:} a sentence of your own with \emph{o artesanato}
  \item \emph{Recall:} say the Portuguese for collection, then the Portuguese for to collect, then say \emph{o artesanato} again
\end{itemize}

\begin{tcolorbox}[breakable,colback=teal!4,colframe=teal!35!black,title={Before you move on}]
What does o artesanato mean? (\textquotedblleft{}Handicrafts, crafts\textquotedblright{}.) Say it once more.
\end{tcolorbox}
\section[o barro]{o barro — clay}

\label{lesson:PT-C504-barro}



\begin{quote}
Before the new one: say the Portuguese for to collect, then the Portuguese for handicrafts.
\end{quote}

\subsection*{You'll want to know: o barro}
\textbf{o barro} — \textquotedblleft{}clay\textquotedblright{}.

\begin{quote}
As crianças fizeram pratos de barro.
\end{quote}

\begin{quote}
\emph{The children made plates out of clay.}
\end{quote}

The famous painted roosters are made of \emph{barro}.

\begin{cousinweb}
\emph{barro} comes from Vulgar Latin \emph{barrum} \textquotedblleft{}clay, mud\textquotedblright{} (EtymDB), probably a word from a language spoken in Iberia before Latin.
\end{cousinweb}

\begin{grammarlens}[title={What you've built}]
One more word for telling, explaining and describing at B1.
\end{grammarlens}

\subsection*{Your turn}
\begin{itemize}
\raggedright
  \item \emph{Say it:} \emph{o barro}
  \item \emph{Say it:} \emph{o barro}, once more
  \item \emph{Say it:} a sentence of your own with \emph{o barro}
  \item \emph{Recall:} say the Portuguese for to collect, then the Portuguese for handicrafts, then say \emph{o barro} again
\end{itemize}

\begin{tcolorbox}[breakable,colback=teal!4,colframe=teal!35!black,title={Before you move on}]
What does o barro mean? (\textquotedblleft{}Clay\textquotedblright{}.) Say it once more.
\end{tcolorbox}
\section[a olaria]{a olaria — pottery}

\label{lesson:PT-C504-olaria}



\begin{quote}
Before the new one: say the Portuguese for handicrafts, then the Portuguese for clay.
\end{quote}

\subsection*{You'll want to know: a olaria}
\textbf{a olaria} — \textquotedblleft{}pottery\textquotedblright{}.

\begin{quote}
Na aldeia ainda há uma olaria antiga.
\end{quote}

\begin{quote}
\emph{There is still an old pottery in the village.}
\end{quote}

The potter is \emph{o oleiro}.

\begin{cousinweb}
\emph{olaria} is built on \emph{oleiro} \textquotedblleft{}potter\textquotedblright{}, which goes back to Latin \emph{olla} \textquotedblleft{}a pot\textquotedblright{} (Lewis \& Short).
\end{cousinweb}

\begin{grammarlens}[title={What you've built}]
One more word for telling, explaining and describing at B1.
\end{grammarlens}

\subsection*{Your turn}
\begin{itemize}
\raggedright
  \item \emph{Say it:} \emph{a olaria}
  \item \emph{Say it:} \emph{a olaria}, once more
  \item \emph{Say it:} a sentence of your own with \emph{a olaria}
  \item \emph{Recall:} say the Portuguese for handicrafts, then the Portuguese for clay, then say \emph{a olaria} again
\end{itemize}

\begin{tcolorbox}[breakable,colback=teal!4,colframe=teal!35!black,title={Before you move on}]
What does a olaria mean? (\textquotedblleft{}Pottery\textquotedblright{}.) Say it once more.
\end{tcolorbox}
